\subsection{Analysis Procedures}
\label{sec:analysis-procedures}

An analysis procedure determines which finite-element problem FEMaster solves. The same mesh, materials, Sections, loads, and constraints can be reused for different procedures because the governing equations differ even when the physical model is unchanged.

FEMaster supports both \cmd{STEP} with a procedure keyword and
\cmd{LOADCASE, TYPE=...} with FEMaster-specific control cards.
Each procedure entry shows its own definition first and then a
corresponding alternative where appropriate. Both can be used
within the unified input language. Constraint handling, nonlinear
convergence, damping, time stepping, and eigensolver controls are
described in Section~\ref{sec:numerical-controls}.

\paragraph*{Analysis families}

The principal equations are:

\textbf{Linear static}
\[
K u=f,
\]
with the prescribed supports and kinematic constraints applied to the displacement vector \(u\).

\textbf{Nonlinear static}
\[
r(u,\lambda)=f_\mathrm{int}(u)-\lambda f_\mathrm{ext}=0,
\]
solved incrementally with Newton iterations. The load factor \(\lambda\) is prescribed by load control or solved together with \(u\) by arc-length control.

\textbf{Eigenfrequency}
\[
K\phi_i=\omega_i^2 M\phi_i,
\qquad
f_i=\frac{\omega_i}{2\pi}.
\]

\textbf{Linear buckling}
\[
\left(K+\lambda_i K_G\right)\phi_i=0,
\]
where \(K_G\) is the geometric stiffness generated by the reference preload state.

\textbf{Linear transient dynamics}
\[
M\ddot u+C\dot u+Ku=f(t),
\]
integrated in time with the implicit Newmark method.

\textbf{Linear harmonic response}
\[
\left(K-\omega^2M+i\omega C\right)\hat u(\omega)=\hat f(\omega),
\]
solved directly for each requested excitation frequency.

\keywordnestedpage{LOADCASE}

Native \cmd{LOADCASE} begins an analysis definition. \key{TYPE} selects one of the supported procedures:
\begin{itemize}
  \item \val{LINEARSTATIC};
  \item \val{NONLINEARSTATIC};
  \item \val{LINEARBUCKLING};
  \item \val{LINEARSTATICTOPO};
  \item \val{EIGENFREQ};
  \item \val{LINEARTRANSIENT};
  \item \val{LINEARHARMONIC}.
\end{itemize}
\key{NAME} is optional and gives the analysis a descriptive identifier.

Commands inside the loadcase select support and load collectors and configure the numerical procedure. Their validity depends on the analysis type: for example, \cmd{TIME} and \cmd{NEWMARK} belong to \val{LINEARTRANSIENT}, \cmd{FREQUENCIES} belongs to \val{LINEARHARMONIC}, and \cmd{NONLINEAR} belongs to \val{NONLINEARSTATIC}.

A deck may contain several native loadcases, which FEMaster executes in input order. Persistent direct structural conditions (such as Abaqus-style boundary and single-DOF concentrated-load definitions) carry their target values into subsequent analyses until modified. Unnamed native vector loads and supports participate in their current analysis, while named load and support collectors are activated for each analysis by its own \cmd{LOADS} and \cmd{SUPPORTS} selections. The analysis procedure determines which model and solution state it uses.

A loadcase can activate reusable named collectors with \cmd{SUPPORTS} and \cmd{LOADS}. The active direct condition history also participates in the analysis, independently of the collector selections.

\begin{keywordparameters}
\key{TYPE} & required & Supported native analysis type listed above. \\
\key{NAME} & optional & Descriptive loadcase name. \\
\end{keywordparameters}

The following static loadcase assumes that the support collector \val{BC} and load collector \val{SERVICE} were defined before the analysis. Both are activated explicitly; other collectors in the model remain inactive.

\exampleheading

\begin{inputdeck}
*LOADCASE, TYPE=LINEARSTATIC,
 NAME=SERVICE_LOAD
*SUPPORTS
BC
*LOADS
SERVICE
*END
\end{inputdeck}

\paragraph*{STEP-based analysis}
\begin{inputdeck}
*STEP, NAME=SERVICE_LOAD
*STATIC
*BOUNDARY
ROOT, 1, 3, 0.
*CLOAD
TIP, 3, -1000.
*END STEP
\end{inputdeck}

\keywordnestedpage{SUPPORTS}

\cmd{SUPPORTS} selects one or more named support collectors for the active analysis. Their support equations are combined with active direct supports and model constraints such as couplings, ties, connectors, and RBM constraints before solving.

Several collectors may be listed together. This supports modular models, for example one collector for permanent fixture constraints and another for a symmetry condition used only in a particular analysis.

The command is supported by the structural analysis families that use nodal constraints, including static, buckling, eigenfrequency, transient, harmonic, and topology-static analyses.

Abaqus input has no separate collector-selection keyword. Its \cmd{BOUNDARY} cards define the support conditions associated with the Step directly.

The selected support collectors have been defined earlier with \cmd{SUPPORT}. Multiple collector names may be activated in one analysis. Previously inherited direct \cmd{BOUNDARY} prescriptions also remain in the condition history.

\begin{keyworddata}
\val{collector_1, collector_2, ...} & Names of previously defined support collectors. \\
\end{keyworddata}

\exampleheading

\begin{inputdeck}
*LOADCASE, TYPE=EIGENFREQ
*SUPPORTS
FIXED_SUPPORTS, SYMMETRY_SUPPORTS
*NUMEIGENVALUES
10
*END
\end{inputdeck}

Direct boundary conditions can instead be specified with
\cmd{BOUNDARY} in a \cmd{STEP}. This does not activate a named
support collector; see the respective keyword pages.

\keywordnestedpage{LOADS}

\cmd{LOADS} activates one or more named load collectors for the current analysis. The assembled external load combines active direct load history with the loads in these selected collectors and any unnamed loads specified in the current analysis.

For a linear static analysis,
\[
f=f_\mathrm{direct}+\sum_c f_c,
\]
where \(f_\mathrm{direct}\) includes active direct loads and \(f_c\) is the contribution of a selected collector. Unnamed native loads specified inside the active analysis also contribute to the current load vector. For transient or harmonic analysis, amplitudes may make the corresponding load contribution depend on time or excitation frequency according to the supported load definition.

\cmd{LOADS} is used by linear static, nonlinear static, buckling, transient, and harmonic analyses. A free-vibration eigenfrequency analysis has no applied external-force vector and therefore does not use a load collector.

Abaqus Step-local load cards such as \cmd{CLOAD}, \cmd{DLOAD}, and \cmd{DSLOAD} provide the corresponding Step loading directly.

The loadcase activates three previously defined collectors. This separates reusable load definitions from the choice of load combinations in each analysis.

\begin{keyworddata}
\val{collector_1, collector_2, ...} & Names of previously defined load collectors. \\
\end{keyworddata}

\exampleheading

\begin{inputdeck}
*LOADCASE, TYPE=LINEARSTATIC
*LOADS
GRAVITY, PRESSURE, ACTUATOR
*END
\end{inputdeck}

Direct \cmd{CLOAD}, \cmd{DLOAD}, and \cmd{DSLOAD} definitions
may also be specified inside \cmd{STEP}, without first activating
named load collectors.

\keywordnestedpage{END}

Native \cmd{END} closes the current \cmd{LOADCASE}. It is distinct from \cmd{END PART}, \cmd{END INSTANCE}, \cmd{END ASSEMBLY}, and Abaqus \cmd{END STEP}.

The closing keyword ends the active \cmd{LOADCASE}. It does not terminate the entire input file and is not used to close an Abaqus-compatible \cmd{STEP}.

\exampleheading

\begin{inputdeck}
*LOADCASE, TYPE=LINEARSTATIC
*SUPPORTS
BC
*LOADS
SERVICE
*END
\end{inputdeck}

A \cmd{STEP} is instead closed by \cmd{END STEP}.
The two closing keywords apply to different analysis scopes.

\keywordnestedpage{STEP}

\cmd{STEP} begins an Abaqus-style analysis step. Multiple steps can appear in one deck and are executed consecutively in input order when their respective \cmd{END STEP} is reached. Each step contains exactly one procedure keyword: \cmd{STATIC}, \cmd{FREQUENCY}, \cmd{BUCKLE}, \cmd{DYNAMIC}, or \cmd{STEADY STATE DYNAMICS}. The keyword selects the analysis class and the active procedure controls.

\key{NAME} assigns an optional step name. \key{NLGEOM} activates the geometrically nonlinear static procedure when present without a value or set to \val{YES}; \val{NO} and \val{OFF} use linear geometry. \key{INC} sets the positive maximum number of nonlinear increments (default 100). \key{PERTURBATION} selects a linear perturbation step. Loads use a ramped step progression by default; a named \cmd{AMPLITUDE} may be referenced on supported load cards to prescribe their history.

Abaqus-style \cmd{BOUNDARY}, single-DOF \cmd{CLOAD}, labeled \cmd{DLOAD} and \cmd{DSLOAD} rows update persistent condition histories. Subsequent steps inherit their target values until matching entries are changed with \key{OP=MOD} or the family is reset with \key{OP=NEW}. Unnamed native six-component \cmd{CLOAD} vectors and native \cmd{SUPPORT} vectors instead apply to the active analysis. Explicit \cmd{LOADS} and \cmd{SUPPORTS} selections activate named collectors for the current analysis only. The \cmd{SOLVER} extension may be placed after the procedure card. \cmd{END STEP} executes the current analysis and closes its scope.

\begin{keywordparameters}
\key{NAME} & optional & Descriptive Step name. \\
\key{NLGEOM} & optional & Bare/\val{YES} enables geometric nonlinearity; \val{NO}/\val{OFF} disables it. \\
\key{INC} & optional & Positive maximum increment count; default 100. \\
\key{RESUME} & optional & \val{YES} (default) retains converged geometry, material and temperature; \val{NO} restores reference state. \\
\key{PERTURBATION} & optional flag & Selects linear perturbation semantics. \\
\end{keywordparameters}

Here a nonlinear static step declares the procedure first, then supplies its boundary conditions and loads, and finally terminates with \cmd{END STEP}.

\exampleheading

\begin{inputdeck}
*STEP, NAME=LARGE_DEFLECTION,
 NLGEOM=YES, INC=100
*STATIC
0.05, 1.0, 1e-5, 0.1
*BOUNDARY
ROOT, 1, 3, 0.
*CLOAD
TIP, 3, -1000.
*END STEP
\end{inputdeck}

\paragraph*{Sequential steps and condition history}
The following two steps use the same node sets \val{ROOT} and \val{TIP}.
The second step retains the support from the first and modifies the
previous \(F_z\) load on \val{TIP} from \(-100\) to \(-200\).
Each step executes at its own \cmd{END STEP}.
\begin{inputdeck}
*STEP, NAME=FIRST
*STATIC
*BOUNDARY
ROOT, 1, 3, 0.
*CLOAD
TIP, 3, -100.
*END STEP
*STEP, NAME=SECOND
*STATIC
*CLOAD, OP=MOD
TIP, 3, -200.
*END STEP
\end{inputdeck}

\paragraph*{Step selecting reusable collectors}
The same \val{BC} and \val{SERVICE} collectors used in native loadcases
may also be selected from a \cmd{STEP} after the procedure card. The
collector contents are defined before the analysis with \cmd{SUPPORT}
using \key{SUPPORT_COLLECTOR=BC} and \cmd{CLOAD} using
\key{LOAD_COLLECTOR=SERVICE}.
\begin{inputdeck}
*STEP, NAME=SERVICE_FROM_COLLECTORS
*STATIC
*SUPPORTS
BC
*LOADS
SERVICE
*END STEP
\end{inputdeck}

\paragraph*{LOADCASE-based analysis}
\begin{inputdeck}
*LOADCASE, TYPE=NONLINEARSTATIC
*SUPPORTS
BC
*LOADS
SERVICE
*NONLINEAR, CONTROL=LOAD,
 MAX_INCREMENTS=100
*END
\end{inputdeck}

\keywordnestedpage{STATIC}

Abaqus \cmd{STATIC} selects static equilibrium. The supported mapping depends on the Step options.

With neither geometric nonlinearity nor Riks path following, the analysis is linear static:
\[
Ku=f.
\]
A \key{PERTURBATION} Step also uses the linear static formulation.
The structural \cmd{TEMPERATURE} card specifies a nodal target field.
General linear static steps commit this field on success; a perturbation
uses the target only for the incremental calculation and discards it.

If \key{NLGEOM} is enabled or \key{RIKS} is present, the procedure uses nonlinear static equilibrium. The nonlinear residual is
\[
r(u,\lambda)=f_\mathrm{int}(u)-\lambda f_\mathrm{ext},
\]
and Newton iterations solve the linearized equilibrium equation
\[
K_T\Delta u=-r.
\]

Without \key{RIKS}, \(\lambda\) is advanced by load control. In a nonlinear
structural step, any prescribed target temperature acts fully in every
increment and Newton iteration; unlike the mechanical loads, it is not
ramped with \(\lambda\). Temperature changes are not supported with
\key{RIKS} in this implementation. \key{RIKS} activates arc-length path following, in which displacement and load factor are solved together under an additional path constraint. This allows equilibrium paths to continue through limit points where simple monotonic load control may fail.

\key{DIRECT} requests fixed nonlinear load increments rather than adaptive stepping where supported.

For ordinary nonlinear static input, the data positions are initial increment, Step period, minimum increment, and maximum increment. The values are converted to the normalized load path represented by the FEMaster nonlinear analysis. For \key{RIKS}, the first four entries analogously define initial, period, minimum, and maximum arc-length scales. Abaqus Riks termination controls in the later data positions are currently not supported.

\begin{keywordparameters}
\key{RIKS} & optional flag & Select arc-length path following. \\
\key{DIRECT} & optional flag & Request fixed nonlinear load increments. \\
\end{keywordparameters}

\begin{keyworddata}
Ordinary static & Optional \val{initial, period, minimum, maximum}. \\
Riks & First four values define the supported path-increment controls; unsupported later termination fields must be omitted. \\
\end{keyworddata}

\paragraph*{Linear static analysis}

\begin{inputdeck}
*STEP, NAME=SERVICE
*STATIC
*BOUNDARY
ROOT, 1, 6, 0.
*CLOAD
TIP, 3, -1000.
*END STEP
\end{inputdeck}

Ordinary linear static analysis always uses the reference displacement
\(u_0=0\) and reference temperature \(T_{\mathrm{ref}}\) as its
linearization point. Each step solves the total linear problem
\[
K_0\,u = f_{\mathrm{ext}}-f_{\mathrm{int}}(0,T_{\mathrm{target}}),
\qquad K_0=K_T(0,T_{\mathrm{ref}}).
\]
The previously calculated linear displacement is not added to the next
solution: the next linear step recomputes the total displacement from
its current total loads and prescribed boundary conditions. Stress,
strain, internal forces, and section resultants use that same reference
linearization. No historical stress or force fields are required.

With \key{RESUME=NO}, the reference temperature and geometry are restored
before the independent solve. With \key{RESUME=YES}, load definitions and
the latest prescribed temperature may be carried over, but a regular
linear step still solves from reference geometry. A regular linear
step cannot resume a nonlinear static solution; to analyze its
prestressed tangent response, use \key{PERTURBATION}, or start an
independent linear analysis with \key{RESUME=NO}.

\paragraph*{Linear static perturbation}

A resumed \cmd{STEP} with \key{PERTURBATION} instead solves the
incremental tangent problem at the committed base state \((u_0,T_0)\):
\[
K_T(u_0,T_0)\,\Delta u =
 (f_{\mathrm{ext,end}}-f_{\mathrm{ext,start}})
 - \bigl[f_{\mathrm{int}}(u_0,T_{\mathrm{target}})
         -f_{\mathrm{int}}(u_0,T_0)\bigr].
\]
The external start load is the last committed applied load, retained
even when input load definitions change or a preceding nonlinear step
converged at a non-unit load factor. Thus load cards retain their
normal target-value meaning:
a \cmd{CLOAD} of 11000 following a step with 10000 denotes a new total target
of 11000, not an additional 11000. The same assembly applies to all supported
structural load families; load-dependent follower tangents are not included.

The output displacement, reaction force, stress, strain, and section force
fields represent incremental responses. The reported temperature is the
absolute target field. No total displacement, temperature or material
state is committed by the perturbation step. On completion, the parser also
restores the original load and support definitions, including values modified
through \key{OP=MOD}/\key{OP=NEW}; later steps resume the previous equilibrium
with its original loads.

\begin{inputdeck}
*STEP, NAME=PRELOAD, NLGEOM=YES
*STATIC
*CLOAD
TIP, 1, 10000.
*END STEP

*STEP, NAME=PERTURB, PERTURBATION, RESUME=YES
*STATIC
*CLOAD
TIP, 1, 11000.
*NODE OUTPUT
U, RF
*ELEMENT OUTPUT
S, E, SF
*END STEP
\end{inputdeck}

The perturbation procedure requires \key{RESUME=YES}. Contact linearization,
free-body rebalancing and inertia relief are not supported in this mode.

\paragraph*{Nonlinear arc-length analysis}
\begin{inputdeck}
*STEP, NAME=POSTBUCKLING, NLGEOM=YES, INC=200
*STATIC, RIKS
0.05, 1.0, 1e-5, 0.1
*END STEP
\end{inputdeck}

\paragraph*{FEMaster nonlinear loadcase}
\begin{inputdeck}
*LOADCASE, TYPE=NONLINEARSTATIC
*NONLINEAR, CONTROL=ARC_LENGTH, INITIAL_INCREMENT=0.05, MAX_INCREMENTS=200
*END
\end{inputdeck}

\keywordnestedpage{FREQUENCY}

Abaqus \cmd{FREQUENCY} selects free-vibration eigenfrequency extraction. The constrained generalized eigenproblem is
\[
K\phi_i=\lambda_i M\phi_i,
\qquad
\lambda_i=\omega_i^2,
\qquad
f_i=\frac{\sqrt{\lambda_i}}{2\pi}.
\]
The data line begins with the positive number of requested modes.

The mode shapes \(\phi_i\) describe relative deformation patterns; their arbitrary scale is not a physical vibration amplitude. Frequencies are determined by both stiffness and mass, so density, concentrated masses, Sections, and support conditions must all be physically meaningful.

\key{EIGENSOLVER} accepts \val{LANCZOS} or \val{SUBSPACE} as supported Abaqus syntax. FEMaster's own eigensolution settings determine the numerical backend used for the calculation.

\begin{keyworddata}
First value & Positive number of requested eigenpairs. \\
Remaining accepted positions & Additional Abaqus spectral controls not used by the currently supported mapping. \\
\end{keyworddata}

The eigenfrequency step below requests 12 modes of the constrained model. The node set \val{FIXED} must already exist, and meaningful frequencies require stiffness and mass definitions.

\begin{keywordparameters}
\key{EIGENSOLVER} & optional & \val{LANCZOS} or \val{SUBSPACE}; the solver backend is selected internally. \\
\end{keywordparameters}

\exampleheading

\begin{inputdeck}
*STEP, PERTURBATION
*FREQUENCY, EIGENSOLVER=LANCZOS
12
*BOUNDARY
FIXED, 1, 6, 0.
*END STEP
\end{inputdeck}

\paragraph*{Eigenfrequency loadcase}
\begin{inputdeck}
*LOADCASE, TYPE=EIGENFREQ
*SUPPORTS
BC
*NUMEIGENVALUES
12
*END
\end{inputdeck}

\keywordnestedpage{BUCKLE}

Abaqus \cmd{BUCKLE} selects linear eigenvalue buckling. The analysis first establishes the reference stress state caused by the selected preload. That stress state produces a geometric stiffness matrix \(K_G\). Buckling factors then follow from
\[
\left(K+\lambda_iK_G\right)\phi_i=0.
\]
Depending on the sign convention used for \(K_G\), the equivalent generalized eigenproblem may be written with \(-K_G\) on the right-hand side.

The eigenvalue \(\lambda_i\) is a multiplier on the reference preload, not an absolute force. If the preload corresponds to 10 kN and the first factor is 3.2, the linearized critical load associated with that mode is 32 kN. This interpretation is valid only within the assumptions of linear eigenvalue buckling; imperfections, material nonlinearity, contact changes, and nonlinear prebuckling deformation are not represented by the eigenvalue alone.

The first data value is the positive number of requested buckling modes. \key{EIGENSOLVER} accepts \val{LANCZOS} or \val{SUBSPACE} as supported Abaqus syntax.

The buckling step first forms the reference stress state from its applied load and then computes load factors. The selected nodes, load and element properties must exist before this analysis.

\begin{keywordparameters}
\key{EIGENSOLVER} & optional & \val{LANCZOS} or \val{SUBSPACE}; the solver backend is selected internally. \\
\end{keywordparameters}
\begin{keyworddata}
\val{count} & Number of requested buckling modes. \\
\end{keyworddata}

\exampleheading

\begin{inputdeck}
*STEP, PERTURBATION
*BUCKLE
5
*BOUNDARY
ROOT, 1, 3, 0.
*CLOAD
TIP, 1, -1000.
*END STEP
\end{inputdeck}

\paragraph*{Buckling loadcase}
\begin{inputdeck}
*LOADCASE, TYPE=LINEARBUCKLING
*SUPPORTS
BC
*LOADS
PRELOAD
*NUMEIGENVALUES
5
*SIGMA
0.0
*END
\end{inputdeck}

\keywordnestedpage{DYNAMIC}

Abaqus \cmd{DYNAMIC, DIRECT} selects FEMaster's linear implicit transient analysis. The structural equation is
\[
M\ddot u(t)+C\dot u(t)+Ku(t)=f(t).
\]
The current supported mapping uses a fixed time increment and therefore requires \key{DIRECT}. Geometric nonlinearity is not supported by this linear transient procedure.

The first two data values are the fixed time increment \(\Delta t\) and the positive Step period \(T\). The analysis runs from \(t=0\) to \(t=T\) with the specified increment. Loads may vary with time through amplitudes.

Native FEMaster exposes the same analysis through \val{LINEARTRANSIENT}, \cmd{TIME}, \cmd{NEWMARK}, optional \cmd{DAMPING}, and output-cadence controls.

The transient step integrates the response for one second using 0.01-second increments. For time-dependent loading, define an \cmd{AMPLITUDE} and reference it from \cmd{CLOAD}.

\begin{keywordparameters}
\key{DIRECT} & required flag & Fixed-increment implicit time integration. \\
\end{keywordparameters}
\begin{keyworddata}
\val{dt, period} & Time increment and positive step duration. \\
\end{keyworddata}

\exampleheading

\begin{inputdeck}
*STEP
*DYNAMIC, DIRECT
0.01, 1.0
*BOUNDARY
ROOT, 1, 3, 0.
*CLOAD, AMPLITUDE=RAMP
TIP, 1, 100.
*END STEP
\end{inputdeck}

\paragraph*{Transient loadcase}
\begin{inputdeck}
*LOADCASE, TYPE=LINEARTRANSIENT
*SUPPORTS
BC
*LOADS
DYNAMIC_LOAD
*TIME
0., 1., 0.01
*NEWMARK
0.25, 0.5
*END
\end{inputdeck}

\keywordnestedpage{STEADY STATE DYNAMICS}

Abaqus \cmd{STEADY STATE DYNAMICS, DIRECT} selects direct linear harmonic response. For harmonic excitation at angular frequency \(\omega=2\pi f\), FEMaster solves the complex system
\[
\left(K-\omega^2M+i\omega C\right)\hat u=\hat f.
\]
The result \(\hat u\) contains amplitude and phase information. Resonances appear when the excitation frequency approaches a natural frequency and damping is small.

\key{DIRECT} is required. \key{INTERVAL} currently supports \val{RANGE}. \key{FREQUENCYSCALE} may be \val{LINEAR} or \val{LOGARITHMIC}.

Each data record contains lower frequency, optional upper frequency, optional point count, bias, and scale factor. If the upper frequency is omitted or zero, only the lower frequency is solved. For a range, the upper frequency must be at least as large as the lower. The default point count is 20. Bias values other than 1 and scale factors other than 1 are not supported. A logarithmic range requires a positive lower frequency.

Native FEMaster specifies a single equidistant frequency sweep per \cmd{FREQUENCIES, SCALE=LINEAR} card using the lower bound, upper bound and point count. See Section~\ref{sec:numerical-controls} for the spacing formula and the corresponding native loadcase example.

This step performs a frequency-response analysis over 10--1000 Hz. The harmonic force is applied to the previously defined node set \val{TIP}; the stiffness and mass properties must be assigned to the model.

\begin{keywordparameters}
\key{DIRECT} & required flag & Solve the direct complex harmonic system. \\
\key{INTERVAL} & optional & \val{RANGE} (only supported setting). \\
\key{FREQUENCYSCALE} & optional & \val{LINEAR} or \val{LOGARITHMIC}. \\
\end{keywordparameters}
\begin{keyworddata}
\val{f_min, f_max, count, bias, scale} & Frequency interval, sample count and two unit-valued controls. \\
\end{keyworddata}

\exampleheading

\begin{inputdeck}
*STEP
*STEADY STATE DYNAMICS, DIRECT,
 FREQUENCYSCALE=LINEAR
10., 1000., 100, 1., 1.
*CLOAD, REAL
TIP, 3, -100.
*END STEP
\end{inputdeck}

\paragraph*{Harmonic loadcase}
\begin{inputdeck}
*LOADCASE, TYPE=LINEARHARMONIC
*SUPPORTS
BC
*LOADS
HARMONIC_FORCE
*FREQUENCIES, SCALE=LINEAR
10., 1000., 100
*END
\end{inputdeck}

\keywordnestedpage{END STEP}

\cmd{END STEP} closes the active Abaqus Step. A supported procedure must have been selected within the Step.

This closing keyword executes and finishes the selected \cmd{STEP}. A FEMaster \cmd{LOADCASE} uses \cmd{END} instead.

\exampleheading

\begin{inputdeck}
*STEP
*STATIC
*BOUNDARY
ROOT, 1, 3, 0.
*CLOAD
TIP, 3, -1000.
*END STEP
\end{inputdeck}

The corresponding \cmd{LOADCASE} scope is closed by \cmd{END},
not by \cmd{END STEP}.

